ACTS is a mature tracking software and is already adopted by a number of experiments. The project started out with ambitious goals, and they continue to create new challenges. The major goals are experiment independence and the use of modern C++. There is always another new experiment with a different environment that will require new features and improvements to ACTS. Some of the current challenges are: accommodating experiments with telescope detectors, providing computationally efficient software with optimal tracking performance at low momentum and in high multiplicity environments, and making use of time information. Concurrently, the C++ standard is constantly evolving and new features can be leveraged to improve ACTS. Before ACTS, people had to write their own tracking software from scratch or extract it from other experiments. This requires personnel, is time-consuming and error-prone, and depends on existing experience or knowledge transfer. ACTS aims to fill this gap by being general purpose, scalable to support future experiments, and by providing a platform to exchange tracking challenges and solutions.

The developments part of this thesis presents the work performed on various parts of ACTS. All the developments presented in this thesis were done in close collaboration with the ACTS community and have been merged into the mainline of ACTS. General improvements to ACTS which are not specific to any experiment or use case are described in Chapter \ref{ch:dev/general}. The developments on the stepper and navigator are discussed in detail in Chapter \ref{ch:dev/stepper} and \ref{ch:dev/navigator}, respectively. These were necessary to support one of the main focuses of this thesis: the developments of the track finding component, which is described in Chapter \ref{ch:dev/track_finding}. These developments were driven by the need to study and improve the track finding performance in ACTS. Chapter \ref{ch:dev/dense} and \ref{ch:dev/4d_tracking} present the developments on tracking in dense material and 4D tracking, respectively. Both add new capabilities to ACTS, supporting a wider range of experiments. This part of the thesis concludes with a summary of what was discussed in the previous chapters and an outlook on future developments in Chapter \ref{ch:dev/summary}.
